%------------------------------------------------------------------------------%
\begin{question}
  Avalie
  \(
    \lim_{x \to \infty} \left(\sqrt{x^2+3x} - x\right)
  \)
\end{question}

%------------------------------------------------------------------------------%
\begin{resolution}{Solução}
  Não podemos substituir diretamente, mas sabemos que
  \[
    \lim_{x \to \infty} \sqrt{x^2+3x} = \infty
  \]
  \[
    \lim_{x \to \infty} x = \infty
  \]
  portanto temos uma indeterminação \(\infty-\infty\)
  e não podemos dizer nada sobre o limite
\end{resolution}

%------------------------------------------------------------------------------%
\begin{resolution}{Solução}
\begin{columns}[t]
\column{0.5\linewidth}
  \begin{align*}
    f(x)
    & = \sqrt{x^2+3x} - x \\
    & = \left(\sqrt{x^2+3x} - x\right)\; \frac{\sqrt{x^2+3x} + x}{\sqrt{x^2+3x} + x} \\
    & = \frac{x^2+3x - x^2}{\sqrt{x^2+3x} + x} \\
    & = \frac{3x}{\sqrt{x^2+3x} + x} \\
    & = \frac{3x}{\sqrt{x^2\left(1+\dfrac{3}{x}\right)} + x}
  \end{align*}
\column{0.5\linewidth}
  \begin{align*}
    f(x)
    & = \frac{3x}{x\sqrt{1+\dfrac{3}{x}\;} + x} \\
    & = \frac{3x}{x\left(\sqrt{1+\dfrac{3}{x}\;} + 1\right)} \\
    & = \frac{3}{\sqrt{1+\dfrac{3}{x}\;} + 1}
  \end{align*}
\end{columns}
\end{resolution}

%------------------------------------------------------------------------------%
\begin{resolution}{Solução}
  \begin{align*}
    \lim_{x \to \infty} f(x)
    & = \lim_{x \to \infty} \frac{3}{\sqrt{1+\dfrac{3}{x}\;} + 1} \\
    & = \frac{3}{\sqrt{\displaystyle 1+\lim_{x \to \infty}\left(\frac{3}{x}\right)\;} + 1} \\
    & = \frac{3}{\sqrt{1+0\;} + 1} \\
    & = \frac{3}{2}
  \end{align*}
\end{resolution}

%------------------------------------------------------------------------------%